\justifying
\NSFCBodyText

本节为公开示例。实际申请书中的成果、项目、能力、资料和授权状态均须由申请人按附件及证明材料据实填写，不应直接沿用示例中的占位内容。

\subsubsection{与本项目相关的研究工作积累}

\subsubsubsection{主题与理论积累}
申请人围绕【研究主题】持续开展【文献研究、案例分析、调查研究或方法应用】等工作，已形成与本项目相关的问题意识和分析基础。代表性成果可按“【论文、著作或项目名称】---【与本项目的关系】”的方式列示，并与申请书附件、成果清单保持一致。上述积累应具体说明其如何支持职业转型事件识别、机制解释或比较检验，而非仅罗列成果数量。

\subsubsubsection{方法与资料能力}
申请人及项目组具备【语种阅读、资料检索、文本或视听材料分析、质性或定量研究】等与研究任务相匹配的能力。已经掌握的方法、软件和工具，应说明其实际用途、版本或授权状态；尚处于学习、试用或拟采购状态的条件，应如实标为“拟补足”，不得表述为既有能力。

\subsubsubsection{可复用资源与边界}
如申请人已完成相关项目或预研究，可据实说明可复用的资料整理经验、编码规则或分析流程，并注明【项目名称、编号、时间范围、授权状态】。涉及受限资料、访谈记录、平台数据或第三方资源时，应说明其合法来源和使用边界；不能公开或尚未获授权的资源不作为本项目按期完成的前提。

\subsubsection{项目组能力与任务分工}

\subsubsubsection{任务匹配}
项目组拟由具有【相关学科背景】的成员承担文献梳理、资料目录建立、编码复核、过程分析和成果整理等任务。实际成员姓名、职称、成果与分工应按申请书要求据实填写，并与参与者简历和依托单位条件一致。

\subsubsubsection{研究质量保障}
项目组将通过资料来源记录、编码复核、阶段性讨论和反例检验保障研究质量。若需要语种、方法或行业咨询，将在依托单位条件、经费和书面确认的合作范围内安排；未落实的支持仅作为拟解决条件，不作为既有基础。

\subsubsection{研究风险的应对措施}

本项目已识别的主要风险及其应对方式如下。每条均按“风险描述---早期信号---应对预案”的顺序说明，预案包含可执行的替代路线与降级目标，确保风险发生时仍有可交付成果。实际申请书中应结合本人资料条件与研究设计据实调整，不宜照搬。

\subsubsubsection{风险一：资料可得性与授权不确定（资源风险）}
\textbf{风险描述：}【受限档案、机构记录、访谈对象或平台数据】的获取依赖第三方授权，存在授权周期长于预期、授权范围窄于申请或部分资料不可用的可能，可能影响样本覆盖面与比较设计的完整性。

\textbf{早期信号：}立项后【3---6 个月】内授权申请未获书面答复；或已获授权资料的实际可用比例低于纳入标准所需规模；或关键子群体的资料缺口无法通过既有渠道补足。

\textbf{应对预案：}第一，在项目启动阶段同步推进两条以上资料渠道，避免单一来源依赖；第二，若受限资料无法按期获得，改以【公开档案、已刊文献、公开访谈记录或二手数据库】替代，并在方法部分显式说明来源变更对结论外推范围的限制；第三，降级目标为缩小比较单元数量、聚焦【核心子群体】完成机制分析，仍可交付资料目录、编码手册与主体机制研究论文，只是比较性结论的适用边界相应收窄。

\subsubsubsection{风险二：编码一致性与解释效度不足（技术与方法风险）}
\textbf{风险描述：}本项目依赖对文本或视听材料的人工编码，存在编码者理解偏差、类目边界模糊、关键节点判定不一致等问题，可能削弱机制解释的可信度。

\textbf{早期信号：}预编码阶段双人独立编码的一致性指标低于预设阈值；或同一材料在复核中出现反复改判；或新增材料频繁触发类目增补，说明编码框架尚未稳定。

\textbf{应对预案：}第一，先在小规模样本上完成编码手册迭代，达到预设一致性阈值后再全面铺开；第二，保留全部分歧记录并由第三方复核仲裁，将争议个案单独归档而非强行归类；第三，若一致性长期不达标，改用【更粗粒度的类目体系】并辅以典型个案深描，降级目标为以少量高信度个案支撑机制假设，仍可交付编码手册、一致性检验报告与个案分析成果。

\subsubsubsection{风险三：关键节点资料整理进度滞后（进度风险）}
\textbf{风险描述：}资料清洗、目录建立与过程追踪工作量易被低估，且与项目组成员的教学科研任务存在时间冲突，可能导致中期任务顺延、压缩后期分析与成果整理时间。

\textbf{早期信号：}第一年度末资料入库与编码完成比例低于年度计划的【70\%】；或阶段性讨论中反复出现同一环节未完成；或单一成员承担的任务出现持续积压。

\textbf{应对预案：}第一，将资料整理拆分为可独立验收的批次，按批次滚动推进，避免整体延期；第二，设置季度检查点，一旦某批次滞后即在项目组内重新分配任务或调整批次顺序，优先保障进入分析环节所必需的核心批次；第三，若进度仍不可控，压缩非核心时段或非核心地区的资料范围，降级目标为完成核心时段的完整分析链条，确保论文与数据成果按期产出，扩展部分顺延至项目结题后继续完善。

\subsubsection{与本项目的衔接}
上述研究积累旨在支撑本项目的资料纳入与编码、关键节点过程追踪以及比较性边界检验。申请人应在正式文本中逐项说明“已有何种基础---支持哪一项任务---仍缺何种条件”，使研究基础与研究内容、工作条件和年度计划保持一致；风险应对部分则应与年度研究计划的里程碑逐年对齐，说明各年度一旦触发早期信号时的降级路径与仍可交付的成果。
